%% preamble/packages.tex — paper-presentation document-level packages.
%% The custom SUST theme lives in preamble/sustbeamer.sty and is loaded
%% from main.tex BEFORE this file (it must precede graphicx usage that
%% references its color definitions).

% --- Graphics & figures ---
\usepackage{graphicx}
\graphicspath{{figures/}}

% --- Math, units ---
\usepackage{amsmath, amssymb}
\usepackage{siunitx}
\sisetup{scientific-notation=false, round-mode=places, round-precision=3}
\DeclareSIUnit\angstrom{\text{\AA}}

% --- Tables ---
\usepackage{booktabs}
\usepackage{array}
\usepackage{multirow}

% --- Code listings (lighter than the thesis templates; presentations need
%     to stay glance-friendly) ---
\usepackage{listings}
\usepackage{xcolor}
\lstset{
  basicstyle=\ttfamily\footnotesize,
  keywordstyle=\color{sustBlue}\bfseries,
  commentstyle=\color{sustGray}\itshape,
  stringstyle=\color{sustAccent},
  backgroundcolor=\color{sustSoft},
  frame=leftline, framerule=1.5pt, rulecolor=\color{sustBlue},
  xleftmargin=8pt, framexleftmargin=5pt,
  breaklines=true, columns=flexible
}

% --- TikZ for inline schematics & flow diagrams ---
\usepackage{tikz}
\usepackage{pgfplots}
\pgfplotsset{compat=1.18}
\usetikzlibrary{arrows.meta, positioning, calc, shapes.geometric, fit, backgrounds}

% --- Hyperlinks & PDF metadata ---
\usepackage{hyperref}
\hypersetup{
  colorlinks=true,
  linkcolor=sustBlue, citecolor=sustGray, urlcolor=sustBlueLight,
  pdftitle={Title of Your Presentation},
  pdfauthor={Author Name},
  pdfsubject={B.Sc. / M.S. Project Defense}
}

% --- Bibliography helper (we use natbib semantics for \cite / \citep)
\usepackage[numbers,sort&compress]{natbib}

% --- Custom presentation helpers ---
% \highlight<...>{...}  — call-out box for one phrase (no overlay support)
\newcommand{\highlight}[1]{\colorbox{sustSoft}{\textbf{#1}}}

% \timenote{...} — tiny gray marginalia reminding the presenter of pacing
\newcommand{\timenote}[1]{{\color{sustGray}\scriptsize\textit{\#1}}\par}

% \resultbox{title}{contents} — full-bleed blue-accented result callout
\newcommand{\resultbox}[2]{%
  \begin{beamercolorbox}[wd=.96\paperwidth, sep=8pt]{block body example}
    {\bfseries\color{sustBlueDark}#1}\par
    \medskip
    #2
  \end{beamercolorbox}
}

% \keypoint — an animated incremental bullet sweep (used sparingly)
% Available if the presenter wants to layer disclosure of multiple bullets.
\setbeamercovered{transparent=20}
